\clearpage\begingroup\SourceLeaf{618b}
\ArithmeticHead{63}{R}{XI. SCHLUSSBETRACHTUNG.}
\renewcommand{\thefootnote}{\arabic{footnote})}\setcounter{footnote}{1}
\newcommand{\MSecondsComma}{\mathord{\overset{\scriptscriptstyle\prime\prime}{,}}}
\begin{center}XI.\\[3pt]\textls[120]{Schlussbetrachtung.}\end{center}
\noindent\hspace*{\parindent}Bei den ausgedehnten Zahlenrechnungen, die \textsc{Gauss} von früher Jugend
an bis ins hohe Alter ausführte, und zwar meist ohne fremde Hülfe, konnte
es nicht ohne Fehler abgehen, und wir haben bereits an verschiedenen Stellen
angemerkt, dass \textsc{Gauss} sich verrechnet hat, oder dass er ein Ergebnis teilweise
durchgestrichen und umgeändert hat. Bei der Betrachtung, die wir über diese
Rechenfehler anstellen wollen, müssen wir zwei Abschnitte unterscheiden;
der eine umspannt \textsc{Gauss}’ Jugendzeit, der andere seine späteren Lebensjahre.

In seiner Jugendentwickelung, wo er, wie wir gesehen haben, den Problemen meist rechnerisch zu Leibe ging, wo er sich durch rechnerische Betätigung das Rüstzeug schuf für tiefere Studien sowohl, als auch für die Praxis,
fällt uns seine erstaunliche Sicherheit im Zahlenrechnen auf, die es bewirkt,
dass die Zahl der Fehler, gemessen an der Fülle des bewältigten Materials,
auffallend gering ist. Man könnte daher vermuten, dass er zu jener Zeit sich
durch \textls[120]{Kontrollrechnungen} gesichert habe.

Eine Rechnung aus der Zeit seiner Jugendentwickelung, die man als
Kontrollrechnung ansprechen könnte, haben wir auf S. 15 ff. vorgeführt; doch
bereits bei diesem Beispiel zeigt sich, dass die Kontrolle höchst wahrscheinlich nicht der Hauptzweck der Arbeit war. Wir sahen auch bereits, dass
er z. B. bei seinen Tafeln der quadratischen Reste keine Kontrollrechnungen
ausgeführt hat, auch keine, oder wenigstens keine ausreichenden, in den Tafeln
zur Cyklotechnie, und da er, wie wir noch sehen werden, auch in späteren
Jahren den Kontrollrechnungen aus dem Wege ging\footnote{Es darf nicht unerwähnt bleiben, dass er sowohl die speziellen, als auch die allgemeinen Störungen
der Pallas doppelt gerechnet hat, wohl weniger, um die zweite Rechnung mit genaueren Elementen durchzuführen, als vielmehr, um die Richtigkeit zu prüfen. Es wäre offenbar einfacher gewesen, die Rechnung
einmal unter Prüfung auf Rechenfehler zu machen. Eine solche Prüfung hat er indessen nirgends angestellt.}, so ist anzunehmen,
dass er sie auch in seiner Jugendzeit nicht besonders bevorzugte.
\clearpage\SourceLeaf{619}\ArithmeticHead{64}{L}{PHILIPP MAENNCHEN, GAUSS ALS ZAHLENRECHNER.}
\setcounter{footnote}{0}
\noindent\hspace*{\parindent}Dass er, wenn es darauf ankam, eine bewunderungswürdige Geschicklichkeit besass, Kontrollrechnungen im grössten Stil auszuführen, zeigt seine denkwürdige Kritik von \textsc{Vegas} \textit{Thesaurus Logarithmorum}, wo er in scharfsinniger
Weise feststellt, dass in diesem Werke unter den 68038 Logarithmen nicht
weniger als 47746 ungenaue zu erwarten sind\footnote{Werke III, S. 257—264.}. Auch bei dem schon erwähnten \textit{Calculus numerico-exponentialis}, sowie bei der Aufsuchung des arithmetisch-geometrischen Mittels treten vereinzelt Kontrollrechnungen auf, die
aber gewöhnlich den Hauptzweck haben, verschiedene Methoden bezüglich
ihrer Tragweite und praktischen Brauchbarkeit gegen einander abzuwägen.
Dass er sie im allgemeinen \textls[120]{nicht} anwandte, zeigt sich deutlich darin, dass
gerade hier von F. \textsc{Goldscheider} eine Reihe von fehlerhaften Ergebnissen nachgewiesen worden ist\footnote{\textsc{Goldscheider} hat zur Kontrolle jede Rechnung nach mehreren Methoden ausgeführt.}.

Im grossen und ganzen hat er auch schon in seiner Jugend auf \textls[120]{genaue}
Kontrollen verzichtet, dagegen meistert er die Überschlagsrechnung in ungeahnter Vollendung, wobei ihm ein angeborenes Gefühl für die Wahl der Abrundung der auftretenden Zahlen und seine hochgesteigerte Fähigkeit, die
Zahlen zu individualisieren, in gleicher Weise zu Hilfe kommen.

Es darf auch nicht unerwähnt bleiben, dass er durch gewisse Eigenarten
weit mehr vor Fehlern geschützt war als einer, der nach dem Schema rechnet.
Sehr viele Rechenfehler werden beim Addieren einer grösseren Anzahl von
Summanden gemacht, und zwar wohl deshalb, weil während des Addierens
einer Stellenkolonne keine Ruhepause möglich ist. \textsc{Gauss} verstand es, durch
geeignete Zerlegung und Anordnung die Durchführung so zu gestalten, dass
er nur selten mehr als zwei Zahlen zu addieren hatte. Das zeigen viele der
von uns aus dem Nachlass wiedergegebenen Rechenbeispiele; man vergleiche
auch das oben S. 48 durchgeführte Multiplikationsbeispiel.

Dass er eine derartige Addition, sowie auch alle Subtraktionen, von links
nach rechts ausführt, ist eine Eigentümlichkeit, über die er sich gelegentlich
\clearpage\SourceLeaf{620}\ArithmeticHead{65}{R}{XI. SCHLUSSBETRACHTUNG.}
\setcounter{footnote}{0}
\noindent in einem Brief an \textsc{Schumacher}\footnote{Vom 3. Oktober 1844, Werke XII, S. 38; vergl. den Aufsatz von A. \textsc{Galle}, \textit{Gauss als Zahlenrechner}, Materialien für eine wissenschaftliche Biographie von \textsc{Gauss}, Heft IV, 1918, S. 10, 11.} ausführlich äussert, und die ihn durch die
grössere Bequemlichkeit beim Anschreiben auch vor Rechenfehlern schützt.

Dass ihm endlich seine riesigen Tafeln mancherlei Möglichkeiten an die
Hand gaben, um Rechenergebnisse rasch zu kontrollieren, darf wohl als sicher
angenommen werden.

Dagegen haben sich in seinen späteren Arbeiten, besonders in den geodätischen und ganz besonders in den astronomischen, \textls[120]{sehr viele Rechenfehler} eingeschlichen.

Man findet\footnote{Vergl. \textsc{Galle}, a. a. O., S. 13.} die Abweichungen in den Zahlenrechnungen der \textit{Theoria motus},
Werke VII, S. 281 ff. ausführlich von \textsc{Brendel} angegeben, ebenso sind Werke
IX von \textsc{Krüger} überall in den Bemerkungen die Abweichungen aufgeführt
und in den Rentenrechnungen Werke IV, S. 188 sind Unrichtigkeiten von
\textsc{Schering} bemerkt worden. Auch in der \textit{Bestimmung des Breitenunterschiedes}
(Werke IX, S. 1 ff.) sind bereits zu \textsc{Gauss}’ Lebzeiten einige Versehen entdeckt
worden\footnote{Vergl. dazu die folgenden Briefstellen: \textsc{Gauss} an \textsc{Schumacher}, 4. Mai 1829. ».\ .\ . Er [Dr. \textsc{Schmidt}]
brachte mir darauf Zahlen, die anders waren, wie die meinigen, setzte aber zur Erläuterung hinzu, die
Zahlen, die in der \textit{Best. des Br. Untersch.} gedruckt sind, seien durch einen kleinen Rechnungsfehler etwas
entstellt.\ .\ .\ .«
\par\hspace*{1em}\textsc{Gauss} an \textsc{Schumacher}, 30. April 1830. ».\ .\ . Um Ihr Vertrauen zu \textsc{Schmidts} Rechnung zu vergrössern, bemerke ich, dass er die zwei Hauptelemente der Erddimensionen viermal berechnet hat; aber nur
Einmal hat er wegen Rechnungsfehler von neuem gerechnet. Nämlich
\par\hspace*{2em}1. Zahlen in meiner \textit{Breitenbestimmung etc.}
\par\hspace*{3em}Diese hatten einen Rechnungsfehler enthalten, den er später verbesserte, .\ .\ .«}.

Zur Erklärung möge zunächst der Umstand herangezogen werden, dass
\textsc{Gauss} \textls[120]{ungewöhnlich rasch rechnete}. Dies erhellt\footnote{Vergl. \textsc{Galle}, a. a. O., S. 9.} am besten aus seinem
\textit{Journal über die Rechnungen an den Pallasstörungen} (Werke VII, S. 605 ff.). Die
ersten Elemente der Vesta erhielt er, nach der Entdeckung dieses Planeten,
durch nur 10stündige Arbeit (Werke VI, S. 288).

Ein weiterer Grund ist seine Überhäufung mit Berufsgeschäften und der
Umstand, dass er die Untersuchungen so gross anlegte, dass sie die Leistungsfähigkeit eines Einzelnen, selbst eines \textsc{Gauss}, bei weitem überstiegen. In den
\par\noindent X 2 Abh. 6.\hfill 9
\clearpage\SourceLeaf{621}\ArithmeticHead{66}{L}{PHILIPP MAENNCHEN, GAUSS ALS ZAHLENRECHNER.}
\setcounter{footnote}{0}
\noindent weiter unten, im Anhang wiedergegebenen Auszügen aus seinem Briefwechsel
finden wir charakteristische Belege dazu\footnote{Briefe \textsc{Bessel} an \textsc{Gauss} vom 8. u. 19. September 1805 und \textsc{Gauss} an \textsc{Bessel} vom 7. Oktober 1805,
\textit{Briefwechsel Gauss-Bessel}, 1880, S. 17—21.}.

Neben diesen bisher besprochenen Fehlern fällt eine Eigentümlichkeit
auf, die man bei einem theoretisch und praktisch so überragenden Manne,
der noch dazu mit Arbeit überhäuft ist, für unmöglich halten sollte.

Er rechnet nämlich mit überflüssigen Dezimalen. Insbesondere gibt er
bei kleinen Zahlen im Logarithmus ebenso viele Stellen hinter dem Komma
an wie bei grossen Zahlen, wobei dann allerdings die Fehler in den letzten
Dezimalen nichts ausmachen. Freilich sind hier die Dezimalen nicht immer
bloss in den letzten, sondern auch manchmal in den wichtigen vorderen Stellen
falsch.

Die Werte für die Pallasstörungen (Werke VII, S. 543 und sonst\footnote{Vergl. die Briefe von \textsc{Nicolai} an \textsc{Gauss}, Werke VII, S. 579—586.}) hat
\textsc{Gauss} bis zum Betrage von $0\MSecondsComma1$ berücksichtigt und auf $0\MSecondsComma01$ angegeben. Diese
Genauigkeit ist auch illusorisch, weil die nicht berücksichtigten Störungen
höherer Ordnung wohl bis an die Bogenminute heranreichen können.

Die \textit{Tabulae novae motus parabolici} (Werke VII, S. 357 ff.\footnote{Vergl. die Bemerkungen von \textsc{Brendel} ebenda, S. 374.}) hat \textsc{Gauss} auf
5 Dezimalstellen der Bogensekunde berechnet, obwohl 3 genügt hätten, damit
die Interpolation auf 2 Dezimalstellen richtig wird. Hierbei ist allerdings zu
berücksichtigen, dass, wenn die 3. Dezimale eine 5 ist, zur Entscheidung für
die Abrundung eine weitere Rechnung nötig sein kann. Überhaupt aber
hätte es genügt, die Tafel auf ganze Bogensekunden zu geben.

Auffallend ist auch, dass \textsc{Gauss} die genannte Tafel, ebenso wie eine Reihe
anderer Tafeln, in verhältnismässig weitem Intervall berechnet hat, so dass
dann die Interpolation eigentlich die Hauptarbeit ist. Die eben besprochene
Tafel (Werke VII, S. 357) hat er im Intervall des Arguments von 200 berechnet
und wollte sie bis auf $\relax{}^{1}\!/\!_{1000}$\footnote{Vergl. die Bemerkungen a. a. O.} abdrucken. Welch fürchterliche Interpolation!

In seinen geodätischen Arbeiten sind ähnliche Beobachtungen gemacht
worden. \textsc{Frischauf} äussert sich in einem Briefe an \textsc{Brendel} folgendermassen:
»Die Zahlenrechnungen von \textsc{Gauss} sind nicht sehr sicher, trotz seiner be-
\clearpage\SourceLeaf{622}\ArithmeticHead{67}{R}{XI. SCHLUSSBETRACHTUNG.}
\setcounter{footnote}{0}
\noindent rühmten Rechenkünste. Was den Ansatz überflüssiger Stellen anlangt, so war
mir \textsc{Gauss} hierin immer ein Rätsel. Das Unglaublichste hatte er in der Geodäsie dabei geleistet. So z. B. setzt er\footnote{Werke IV, S. 330, Ostw. Klass. 177, S. 67.} $\log(7)$ und $\log 3\,(7)$ auf 10 Stellen
an, wo 2 genügen würden. Ferner setzt er\footnote{Werke IV, S. 272, Ostw. Klass. 177, S. 15.} den Koeffizienten des Gliedes
mit $\left(\frac{q}{100}\right)^5$ auf 6 Dezimalstellen der Sekunde, dessen grösster Wert für $q=6$
nur $0\MSecondsComma00039$ beträgt, also um 7 Stellen zu viel.«

In den \textit{Elementen der Theorie des Erdmagnetismus} (Werke V, S. 150 ff.)
äussert sich\footnote{Vergl. \textsc{Galle}, a. a. O., S. 13.} \textsc{Gauss} zu der gerade bei dieser Arbeit angewandten übertriebenen
Schärfe: »Für jeden Rechnungskundigen ist die Bemerkung überflüssig, dass
diese Bruchtheile an sich keinen Wert haben, da wir noch weit davon entfernt
sind, nur die ganzen Einer mit Zuverlässigkeit ausmitteln zu können: allein
es ist von Wichtigkeit, dass die Beobachtungen mit einem und demselben
bestimmten System von Elementen scharf verglichen werden, und da war kein
Grund vorhanden, an dem, was die Rechnung ergeben hatte, etwas zu verändern, weil durch Weglassung der Dezimalbrüche für die Bequemlichkeit
der Vergleichsrechnungen garnichts gewonnen sein würde.«

Es ist klar, dass durch solche Umständlichkeit seine ohnehin knapp bemessene Zeit noch mehr in Anspruch genommen wurde, wenn er dies auch
gelegentlich bestreitet\footnote{Er behauptet, bei seiner Art, zu rechnen, mache es wenig aus, eine Reihe von Stellen mehr auszuführen. Er spielt dabei vermutlich auf seine Divisionsmethode an, die wir oben S. 9 f. auseinandergesetzt
haben.}.

Wir haben noch\footnote{Vergl. A. \textsc{Galle}, a. a. O., S. 4, 7.} der Hülfen zu gedenken, die ihm bei seiner umfangreichen Rechenarbeit zu teil geworden sind. Bei den Störungsrechnungen (zu
Störungstafeln der Pallas hielt \textsc{Gauss} 6 bis 8 Rechner für nötig) und bei den
Hülfstafeln zur Berechnung der magnetischen Kräfte hat er die Mitarbeit
verschiedener Astronomen (\textsc{Westphal}, \textsc{Encke}, \textsc{Nicolai}, \textsc{Goldschmidt}) in Anspruch genommen\footnote{Vergl. Werke VII, S. 602 und V, S. 152, 177.}. Auch die Hülfe, die \textsc{Bessel}\footnote{Vergl. die Briefe vom 21. und 29. Dezember 1804, \textit{Briefwechsel Gauss-Bessel}, 1880, S. 1—3.} durch Berechnung von
Sonnenörtern und der Koeffizienten in der Entwickelung des reziproken Wertes
der Entfernung zweier Himmelskörper lieh, kann hier erwähnt werden.
\par\hfill $9^{\ast}$
\clearpage\SourceLeaf{623}\ArithmeticHead{68}{L}{PHILIPP MAENNCHEN, GAUSS ALS ZAHLENRECHNER.}
\setcounter{footnote}{0}
\noindent\hspace*{\parindent}Wie \textsc{Gauss} sich zu dem Gebrauch \textls[120]{mechanischer} Hilfsmittel gestellt
haben würde, lässt sich aus dem wohlwollenden Zeugnis, das er dem Erfinder
des Modells einer Rechenmaschine, Professor \textsc{Schiereck}, ausgestellt hat (Werke
X, 1, S. 6), nicht mit Deutlichkeit entnehmen.

Als ihm die Hülfe des Rechenkünstlers \textsc{Dase} angeboten wurde, lehnte
er entschieden ab: er könne sich bei den vielen und grossen Rechnungen,
die er in seinem Leben ausgeführt habe, kaum eines Falles erinnern, wo die
Hülfe von jemand, der bloss mechanische Rechnungsfertigkeit gehabt hätte,
ihm von irgend einem Nutzen hätte sein können. Auf weiteres Drängen hin
verwies er \textsc{Dase} auf ein Gebiet, wo ausdauernde rechnerische Betätigung noch
eine bemerkenswerte Lücke ausfüllen konnte, nämlich die damals noch nicht
geleistete Feststellung der Primzahlen in der \textls[120]{7., 8. und 9. Million}\footnote{Vergl. die Werke XII, S. 39—48 abgedruckten Briefe sowie den weiter unten im Anhang [4.]
wiedergegebenen Brief von \textsc{Gauss} an \textsc{Encke}.}.

Auch an den \textls[120]{graphischen Methoden} ist er nicht achtlos vorbeigegangen. In einem Brief an \textsc{Gerling}, vom 29. Mai 1851, erwähnt er bei
der Beurteilung des bei ihm als Assistenten eingetretenen \textsc{Klinkerfues} eine
elegante graphische Konstruktion der Lösung der Hauptgleichung bei der Berechnung einer Planetenbahn, die die Form hat
\[
a\mathord{.}\sin^4 z=\sin(z+b).
\]

Zusammenfassend können wir sagen, dass \textsc{Gauss} vermutlich, überhaupt
durch Arbeiten als beobachtender Astronom und als messender Geodät, Arbeiten, die ihm leider nicht von geschulten Hilfskräften abgenommen wurden,
verhindert war, seine gross angelegten Rechnungen mit der nötigen Ruhe
durchzuführen und zu prüfen. \textsc{Gauss} klagt auch öfter darüber, und noch
mehr über den Umstand, dass ihm diese Dinge, sowie auch seine Lehrtätigkeit die beste Zeit und Stimmung zu zusammenhängenden Arbeiten rauben\footnote{\textsc{Gauss} an \textsc{Olbers}, 13. Februar 1821, Auszug im Anhang [5.].}.
— Denn gerade in dem Gebiet, in dem er durch seine \textit{Disquisitiones arithmeticae} neue Bahnen eröffnet hat, trug er sich mit gewaltigen Plänen, über
die er in seinem Briefwechsel hie und da Andeutungen machte. So äusserte
er\footnote{Siehe Werke X, 1, S. 75.} gelegentlich der Aufforderung von \textsc{Olbers}, sich an die Lösung des grossen
\textsc{Fermat}schen Problems zu machen, er hoffe, eine Theorie zu entwickeln, aus
\clearpage\SourceLeaf{624}\ArithmeticHead{69}{R}{XI. SCHLUSSBETRACHTUNG.}
\setcounter{footnote}{0}
\noindent der der \textsc{Fermat}sche Satz als ein höchst unwesentliches Korollar hervorgehen
würde. Aber, wie er selbst in dem erwähnten Briefe an \textsc{Olbers} (Anhang [5.])
ahnungsvoll prophezeit, so ist es gekommen: Manche seiner zahlentheoretischen
Arbeiten sind für uns verloren, und damit ist uns auch die Möglichkeit genommen, zu untersuchen, inwieweit ihre Ergebnisse seine Methoden beim Zahlenrechnen beeinflusst haben.

Nach dem bisher Dargestellten dürfen wir, um noch einmal kurz zusammenzufassen, folgenden Entwickelungsgang des Zahlenrechners \textsc{Gauss} vermuten:

In früher Jugend das Spiel mit dem arithmetisch-geometrischen Mittel,
sowie die Versuche, das Verteilungsgesetz der Primzahlen zu ermitteln. Erfolg: Bedeutende Fertigkeit im Addieren\footnote{Er kann nach seiner eigenen Aussage (Brief an \textsc{Olbers} vom 7. Januar 1815, zitiert bei \textsc{Galle},
a. a. O. S. 12) mit der Summe zweier gegebener Zahlen, ohne diese Summe selbst vor sich zu haben, sogleich in eine Tafel eingehen. Ein einfaches Beispiel haben wir auf S. 14, in der Fussnote angegeben.}, Multiplizieren, Wurzelziehen, Zerlegen grosser Zahlen in Summen von Quadraten und damit auch in Primfaktoren.

Dann folgt die Verwertung der gewonnenen Gewandtheit zur Aufstellung
von Tabellen: Quadratzahlen, Dezimalbruchtabelle, Logarithmen, Anfänge der
Cyklotechnie u. a. m. Erfolg: Noch viel mehr gesteigerte Fertigkeit im Zahlenrechnen, in der Auswertung unendlicher Reihen (Wurzeln, Logarithmen u. a. m.),
Faktorenzerlegung grosser Zahlen, bezw. deren Erkennung als Primzahl; eine
Fülle von zahlentheoretischen Eigenschaften, die Grundlage der \textit{Disquisitiones}.

Die unvergleichliche Fülle von Anregungen, die ihm seit 1795 das Studium
von \textsc{Lagrange}, \textsc{Legendre} und vor allem von \textsc{Euler} bietet, erhöht seine zahlentheoretische Einsicht und seine Rechenfertigkeit in gleicher Weise.

Die hoch gesteigerte Fertigkeit im numerischen Rechnen wird nun auf
die empirische Ermittelung von mathematischen Gesetzen auf den verschiedensten Gebieten mit wunderbarem Erfolg angewandt, daneben aber auch auf
die Praxis als Astronom, Physiker, Geodät.

Die induktiv gefundenen rein mathematischen Ergebnisse werden in harter
Arbeit bewiesen, aber zu dieser Arbeit gehört Zeit, die ihm nach der Übernahme der allzuviel Zeit absorbierenden Ämter nur in beschränktem Masse
zur Verfügung steht.
\clearpage\SourceLeaf{625}\ArithmeticHead{70}{L}{PHILIPP MAENNCHEN, GAUSS ALS ZAHLENRECHNER.}
\setcounter{footnote}{0}
\noindent\hspace*{\parindent}Die Spuren des Gewaltigen im Reiche der Zahlen verlieren sich darum
von hier an immer mehr, und es wird schwer halten, sie weiter zu verfolgen.
Soviel aber darf sicher behauptet werden: Die Wechselwirkung, die seine
Jugendentwicklung bereits kennzeichnet (Beobachtungen beim Zahlenrechnen,
Erforschung der Gründe für die beobachteten Erscheinungen, Gewinnung
zahlentheoretischer Sätze, Anwendung dieser Sätze zur Vereinfachung des Zahlenrechnens, darauf gestützt wieder neue Beobachtungen u. s. f.), bestand auch
weiter bis in sein hohes Alter\footnote{Dass \textsc{Gauss} auch noch im hohen Alter die Rechenfreudigkeit besass, die schon den Knaben und
Jüngling auszeichnete, erkennen wir aus einem erst vor kurzem aufgefundenen Stück aus seinem Nachlass
(Werke XII, S. 10). \textsc{Gauss} hatte sich 1852 gelegentlich einmal mit einem Rätsel beschäftigt, dessen Lösung
zu einer Gleichung 5. Grades führte; doch waren nur die positiven ganzzahligen Lösungen brauchbar. Zwei
Lösungen waren leicht erkennbar, und nach Beseitigung der entsprechenden Wurzelfaktoren blieb die folgende kubische Gleichung
\[
x^3+180x^2+366x+182=0.
\]
\noindent Obwohl diese Gleichung augenscheinlich keine positiven Wurzeln mehr haben kann, hat sie \textsc{Gauss} dennoch — und zwar auf ungemein geschickte Art — gelöst und die Wurzeln auf 7 Dezimalstellen bestimmt.}; nicht nur der Beruf, sondern auch die angeborene Neigung haben ihn stets wieder zum reinen Zahlenrechnen zurückgeführt, von dem er einst ausgegangen war, und dem er so viele seiner schönsten
Entdeckungen verdankte. So hat er immer wieder den Boden berührt und
durch diese Berührung immer wieder neue Kräfte gewonnen.
\clearpage\endgroup

